\ifdefined\gkver\else\def\gkver{student}\fi
\documentclass[\gkver]{gaokaozhenti}
\begin{document}

\chapter{2001年全国理综新课程}
%% paperType: 理综物理部分

\begin{enumerate}
\item
%% number: 20
%% paperName: 2001年全国理综新课程第 20 题 6 分
%% typeId: 1
%% type: 单选题
%% chapter: 力物体的平衡
%% point: 力矩平衡
%% method: 无
%% score: 6
%% degree: 450
%% duplicateId: 0
%% body:
图中是轮船上悬挂救生艇的装置的简化示意图。A、B 是船舷上的固定箍，以 $N_{1}$、$N_{2}$ 分别表示固定箍 A、B 作用于吊杆的水平力的大小，已知救生艇所受的重力 $P = 1500\UN$，$d = 1\Um$，$l = 0.8\Um$。如吊杆的质量忽略不计，则\xzanswer{C}

\onepicture[h=3.4cm]{figs/20.png}

\fourchoices[answer=C]
{$N_{1} = 1200\UN$，$N_{2} = 0$}
{$N_{1} = 0$，$N_{2} = 1200\UN$}
{$N_{1} = 1200\UN$，$N_{2} = 1200\UN$}
{$N_{1} = 750\UN$，$N_{2} = 750\UN$}

%% answer: C
%% memo:
\memoanswer{
隔离吊杆如图所示，由分析受力知，绳向下的拉力 $T$，箍 A 对杆的水平作用力 $N_{1}$ 和竖直作用力 $N_{A}$，箍 B 对杆的水平作用力 $N_{2}$ 和竖直作用力 $N_{B}$，选 A 为转动轴，由力矩平衡可得 $N_{2}\cdot d = T\cdot l$；选 B 为转动轴，由力矩平衡可得 $N_{1}\cdot d = T\cdot l$，联立解得 $N_{1} = 1200\UN$，$N_{2} = 1200\UN$，故 C 正确。

\onepicture[h=3.4cm]{figs/20_an.png}
}

\end{enumerate}

\end{document}
